% SPDX-License-Identifier: Apache-2.0
% covers: Doorbell
\datasheet{Doorbell}{Razboj's doorbell: the register a program writes a display list's count into, and the line that wakes the rasteriser.}

\dsfacts{\code{gpu/razboj/src/doorbell.rs}}{\code{razboj::doorbell::Doorbell}}%
{\code{//docs:vreteno}}

\subsection*{Function}

\code{Doorbell} holds the count \code{Raster} reads at its
\code{CTRL}: how many entries of the display list to draw, zero when
there is nothing to draw. A program writes the list into memory, then
the count; the rasteriser draws the list and writes the count back to
zero. The count is a register rather than a word of memory, so that
it holds zero from reset, and no word a loader left behind can start
a list nobody wrote (issue 985).

It drives \code{ring} high while the count is not zero, and the
rasteriser reads the count only then, so an idle rasteriser puts no
reads on the bus. A second word reads the rasteriser's \code{idle}
line.

It is an AXI-Lite peripheral whose link is one port, \code{bus}, a
\code{LitePort} of the five channels: \code{bus\_aw}, \code{bus\_ar}
and \code{bus\_w} in, \code{bus\_b} and \code{bus\_r} out. Its other
ports are \code{rst} and \code{idle} in, and \code{ring} out.

\subsection*{Parameters}

None.

\subsection*{Register map}

Offsets from the peripheral's base, \code{0x3900} on the flagship, the
peripheral page's tenth slot. The port selects the word by bit 2 of
the address.

\dsregs{Doorbell}

\dstables{Doorbell}

\subsection*{Behaviour}

\begin{itemize}
\item A read is answered on \code{bus\_r} in the cycle its address
beat is taken, with \code{OKAY}. A write needs its address beat, its
data beat and room on \code{bus\_b} in the same cycle, and is answered
at once, with \code{OKAY}.
\item A write of \code{count} sets it to the word's low sixteen bits.
A write of \code{status} is ignored.
\item \code{rst} high sets the count to zero, over a write in the same
cycle.
\item \code{ring} is high while the count is not zero.
\end{itemize}

A program waits for \code{count} to read zero, writes the list, then
writes \code{count}. A write while a list is being drawn starts
nothing: the rasteriser has read the count already, and writes zero
over it at the end.

\subsection*{Verification}

\begin{itemize}
\item \code{//gpu/razboj:razboj\_test} runs
\code{the\_count\_holds\_what\_is\_written\_and\_the\_status\_is\_the\_line}:
the count reads back what is written, \code{ring} follows it, the
status is the line, and a reset empties the count; and
\code{the\_doorbell\_lowers}.
\item \code{//cpu/vreteno:board\_test} runs
\code{razboj\_draws\_a\_list\_rung\_on\_the\_doorbell}: a program on the
board writes a list, rings, waits for the zero and for idle, and reads
the pixels back over the bus.
\end{itemize}

\subsection*{Used by}

\code{vreteno32::board::Board}, on the peripheral page's tenth slot,
beside \code{Raster<32, 2, 10, 480, 0x4200\_0000, 0x4280\_0000,
0x3900>}.

\subsection*{Limits}

The count is sixteen bits. A write's strobes are not looked at.
